\chapter{Khảo sát hiện trạng và các hệ thống liên quan}
\label{ch:khaosat}

\section{Các khái niệm liên quan}
\label{sec:khainiem}

\subsection{Văn phòng chia sẻ}
Văn phòng chia sẻ (co-working space) là không gian làm việc chung có sẵn tiện nghi, được nhiều cá nhân và doanh nghiệp độc lập cùng sử dụng. Mô hình này hình thành ở Mỹ vào giữa những năm 2000 và hiện phổ biến ở nhiều thành phố lớn. Các loại chỗ thường gặp gồm:
\begin{itemize}
    \item \textbf{Chỗ ngồi linh hoạt (hot desk):} khách dùng bất kỳ bàn trống nào trong khu vực chung; giá thấp, phù hợp người làm việc tự do.
    \item \textbf{Chỗ ngồi cố định (dedicated desk):} một bàn dành riêng cho khách trong thời hạn thuê, thường kèm tủ cá nhân.
    \item \textbf{Văn phòng riêng (private office):} phòng kín có khóa, dành cho nhóm hoặc doanh nghiệp nhỏ.
    \item \textbf{Phòng họp (meeting room):} phòng cách âm có màn hình, bảng viết, thường thuê theo giờ.
    \item \textbf{Buồng gọi điện (phone booth) và khu sự kiện (event space):} không gian nhỏ để gọi điện hoặc khu lớn để tổ chức hội thảo.
\end{itemize}
Trong CoSpace, các loại không gian được khai báo trong danh mục \mt{workspace\_types}; dữ liệu mặc định gồm bàn làm việc, phòng họp, văn phòng riêng, buồng gọi điện và khu sự kiện.

\subsection{Đặc điểm vận hành}
Các đặc điểm sau ảnh hưởng trực tiếp đến thiết kế phần mềm quản lý:
\begin{itemize}
    \item \textbf{Tài nguyên hữu hạn theo thời gian:} số chỗ cố định nhưng được cho thuê lại theo từng khoảng thời gian. Một khoảng thời gian bỏ trống là doanh thu mất đi, nên hệ thống phải theo dõi tình trạng từng chỗ theo thời gian.
    \item \textbf{Nhiều đơn vị tính giá:} giá theo giờ, ngày, tuần hoặc tháng; thuê dài thường có đơn giá thấp hơn.
    \item \textbf{Dịch vụ phát sinh:} trong lúc làm việc, khách có thể gọi thêm đồ uống, in ấn, thuê thiết bị. Các khoản này được ghi vào hóa đơn của lượt đặt chỗ và thanh toán trước khi trả chỗ.
    \item \textbf{Tính cộng đồng:} khách muốn gặp đối tác, chuyên gia làm việc trong cùng không gian.
\end{itemize}

\subsection{So sánh với văn phòng truyền thống}
Bảng~\ref{tab:sosanh_vanphong} so sánh hai mô hình theo một số tiêu chí vận hành.

\begin{table}[H]
\centering
\caption{So sánh văn phòng chia sẻ và văn phòng truyền thống}
\label{tab:sosanh_vanphong}
\small
\renewcommand{\arraystretch}{1.2}
\begin{tabularx}{\linewidth}{|P{3.3cm}|L|L|}
\hline
\textbf{Tiêu chí} & \textbf{Văn phòng truyền thống} & \textbf{Văn phòng chia sẻ} \\
\hline
Thời hạn hợp đồng & Dài hạn, thường vài năm & Theo giờ, ngày, tuần, tháng hoặc năm \\
\hline
Chi phí ban đầu & Lớn (thiết kế, nội thất, tiền cọc) & Thấp, chủ yếu là tiền thuê \\
\hline
Điều chỉnh quy mô & Khó mở rộng hoặc thu hẹp nhanh & Tăng giảm số chỗ theo nhu cầu \\
\hline
Chi phí vận hành & Doanh nghiệp tự chi trả (lễ tân, vệ sinh, Internet, điện) & Đã tính trong giá thuê \\
\hline
Tiện ích chung & Tự trang bị & Có sẵn pantry, phòng họp, máy in \\
\hline
Kết nối cộng đồng & Ít giao lưu giữa các doanh nghiệp & Dễ gặp gỡ các thành viên khác \\
\hline
\end{tabularx}
\end{table}

\subsection{Các khái niệm kỹ thuật dùng trong đề tài}
\begin{itemize}
    \item \textbf{Đặt chỗ theo khoảng thời gian:} mỗi đơn có thời điểm bắt đầu và kết thúc, biểu diễn bằng khoảng nửa mở $[start\_at, end\_at)$. Nhờ vậy đơn kết thúc lúc 10:00 không xung đột với đơn bắt đầu lúc 10:00.
    \item \textbf{Xung đột lịch:} hai đơn $A$ và $B$ trên cùng một không gian xung đột khi và chỉ khi
    \begin{equation}
        start\_at_A < end\_at_B \quad \text{và} \quad end\_at_A > start\_at_B
        \label{eq:overlap}
    \end{equation}
    \item \textbf{Chính sách hủy nhiều bậc:} tỷ lệ hoàn tiền phụ thuộc khoảng thời gian từ lúc hủy đến lúc bắt đầu (hoặc từ lúc đặt, với bậc ``thời gian ân hạn''). Khoảng thời gian được tính theo phút và so với các khoảng nửa mở $[min, max)$ của từng bậc.
    \item \textbf{Xử lý lũy đẳng (idempotent):} khi cổng thanh toán gửi lại cùng một thông báo nhiều lần, trạng thái đơn và giao dịch chỉ thay đổi ở lần xử lý đầu tiên.
    \item \textbf{Nhật ký kiểm toán (audit log):} bản ghi lưu người thực hiện, hành động, đối tượng, giá trị cũ, giá trị mới, địa chỉ IP và thời điểm của các thao tác quản trị.
\end{itemize}

\section{Khảo sát các hệ thống hiện có}
\label{sec:khaosat_hethong}

Nhóm khảo sát bốn nhóm giải pháp: các chuỗi co-working trong nước, chuỗi quốc tế có mặt tại Việt Nam, phần mềm quản lý co-working dạng SaaS và cách quản lý thủ công. Các nhận xét dưới đây dựa trên trang web công khai và trải nghiệm đặt chỗ của nhóm tại thời điểm thực hiện đề tài (năm 2026), nên có thể thay đổi theo thời gian.

\subsection{Các chuỗi co-working trong nước}
\begin{itemize}
    \item \textbf{CirCO (TP. Hồ Chí Minh):} chuỗi co-working có nhiều chi nhánh ở các quận trung tâm. Khách xem bảng giá, hình ảnh trên trang web rồi gửi biểu mẫu hoặc gọi điện để được tư vấn. Nhóm chưa thấy sơ đồ chỗ ngồi tương tác, và việc đặt chỗ theo giờ vẫn cần nhân viên xác nhận.
    \item \textbf{Toong:} chuỗi co-working có cơ sở ở Hà Nội và TP. Hồ Chí Minh, chủ yếu bán gói thuê theo tháng, năm và văn phòng ảo qua nhân viên kinh doanh. Nhóm chưa thấy chức năng đặt chỗ theo giờ và thanh toán trực tuyến trên trang web.
    \item \textbf{Dreamplex (TP. Hồ Chí Minh):} tập trung vào khách hàng doanh nghiệp và các sự kiện cộng đồng. Việc đăng ký chủ yếu qua tư vấn viên hoặc biểu mẫu trực tuyến.
    \item \textbf{Các không gian quy mô nhỏ:} nhận đặt chỗ qua Fanpage, Zalo hoặc trực tiếp tại quầy, quản lý bằng bảng tính và đối soát chuyển khoản qua ảnh chụp màn hình.
\end{itemize}

\subsection{Chuỗi quốc tế có mặt tại Việt Nam}
\textbf{Regus / Spaces (tập đoàn IWG)} có cơ sở ở Hà Nội, Đà Nẵng và TP. Hồ Chí Minh. Thành viên có thể đặt phòng họp và mua gói dịch vụ qua ứng dụng và cổng thông tin trực tuyến. Hạn chế đối với nhóm khách hàng mục tiêu của đề tài là giá thuê cao và phương thức thanh toán chủ yếu bằng thẻ quốc tế. Hình~\ref{fig:khaosat_trongnuoc} minh họa giao diện đặt chỗ của CirCO và Regus.

\begin{figure}[H]
  \centering
  \begin{subfigure}[b]{0.48\textwidth}
    \centering
    \includegraphics[width=\textwidth,height=5.2cm,keepaspectratio]{circo-booking-system.png}
    \caption{Giao diện đặt chỗ của CirCO}
    \label{fig:circo_booking}
  \end{subfigure}
  \hfill
  \begin{subfigure}[b]{0.48\textwidth}
    \centering
    \includegraphics[width=\textwidth,height=5.2cm,keepaspectratio]{regus-booking.png}
    \caption{Giao diện đặt phòng họp của Regus}
    \label{fig:regus_booking}
  \end{subfigure}
  \caption{Giao diện đặt chỗ của một số chuỗi văn phòng chia sẻ}
  \label{fig:khaosat_trongnuoc}
\end{figure}

\subsection{Phần mềm quản lý co-working dạng SaaS}
\begin{itemize}
    \item \textbf{Nexudus:} nền tảng quản lý co-working có quản lý thành viên, sơ đồ mặt bằng, thanh toán, ứng dụng di động và tích hợp khóa cửa.
    \item \textbf{OfficeRnD:} nền tảng quản lý không gian làm việc linh hoạt, tích hợp với các bộ công cụ văn phòng và hệ thống CRM.
    \item \textbf{Cobot:} giải pháp gọn nhẹ cho không gian nhỏ, tập trung vào thu phí thành viên và đặt phòng họp.
\end{itemize}
Các nền tảng này có nhiều chức năng nhưng khó áp dụng cho chuỗi co-working nhỏ tại Việt Nam vì ba lý do: chi phí thuê bao tính theo số thành viên hoặc số chỗ, cổng thanh toán tích hợp sẵn chủ yếu là các cổng quốc tế (chưa hỗ trợ VietQR), và quy trình khó tùy biến cho thói quen thanh toán tiền mặt tại quầy.

\subsection{Quản lý thủ công}
Nhiều không gian độc lập dùng kết hợp Google Sheets, sổ ghi chép và nhóm chat để quản lý chỗ trống. Cách này rẻ nhưng dễ ghi đè dữ liệu khi nhiều người cùng sửa, không ngăn được đặt trùng khi có khách đặt đột xuất, khó kiểm soát khách ngồi quá giờ và mất nhiều thời gian để tổng hợp doanh thu cuối tháng.

\section{So sánh các giải pháp}
\label{sec:sosanh}

Bảng~\ref{tab:sosanh_hethong} so sánh các giải pháp theo các tiêu chí liên quan đến bài toán. ``Chưa ghi nhận'' nghĩa là nhóm không tìm thấy chức năng này trong quá trình khảo sát.

\begin{table}[H]
\centering
\caption{So sánh các giải pháp đã khảo sát}
\label{tab:sosanh_hethong}
\footnotesize
\renewcommand{\arraystretch}{1.2}
\begin{tabularx}{\linewidth}{|L|P{2.1cm}|P{1.9cm}|P{1.9cm}|P{1.8cm}|P{1.9cm}|}
\hline
\textbf{Tiêu chí} & \textbf{Chuỗi trong nước} & \textbf{Regus} & \textbf{Nexudus} & \textbf{Thủ công} & \textbf{CoSpace} \\
\hline
Sơ đồ chỗ ngồi tương tác & Chưa ghi nhận & Chưa ghi nhận & Có & Không & Có \\
\hline
Đặt chỗ trực tuyến theo giờ & Một phần & Có & Có & Không & Có \\
\hline
Tự động chống đặt trùng & Thủ công & Có & Có & Thủ công & Có (2 lớp) \\
\hline
Thanh toán VietQR & Chuyển khoản thủ công & Chưa ghi nhận & Chưa ghi nhận & Thủ công & Có (PayOS) \\
\hline
Chính sách hủy tự động & Thủ công & Có & Có & Thủ công & Có, theo chi nhánh \\
\hline
Quản lý nhiều chi nhánh & Một phần & Có & Có & Không & Có \\
\hline
Gợi ý kết nối thành viên & Chưa ghi nhận & Chưa ghi nhận & Chưa ghi nhận & Không & Có \\
\hline
Trợ lý ảo & Chưa ghi nhận & Chưa ghi nhận & Chưa ghi nhận & Không & Có (Gemini) \\
\hline
\end{tabularx}
\end{table}

\noindent\textbf{Nhận xét:} các chuỗi trong nước mới dừng ở mức giới thiệu và nhận yêu cầu đặt chỗ, phụ thuộc nhiều vào nhân viên. Chuỗi quốc tế và phần mềm SaaS có quy trình hoàn chỉnh hơn nhưng chi phí cao và chưa phù hợp với cách thanh toán phổ biến ở Việt Nam. Cách quản lý thủ công rẻ nhưng không mở rộng được khi số chi nhánh tăng.

\section{Khoảng trống và định hướng giải pháp}
\label{sec:gap}

Từ kết quả khảo sát, nhóm xác định sáu khoảng trống và chức năng tương ứng của CoSpace (Bảng~\ref{tab:gap_tinhnang}).

\begin{table}[H]
\centering
\caption{Khoảng trống khảo sát và chức năng tương ứng của CoSpace}
\label{tab:gap_tinhnang}
\small
\renewcommand{\arraystretch}{1.2}
\begin{tabularx}{\linewidth}{|P{0.8cm}|L|L|}
\hline
\textbf{STT} & \textbf{Khoảng trống} & \textbf{Chức năng đề xuất} \\
\hline
1 & Khách khó chọn đúng vị trí mong muốn & Sơ đồ mặt bằng hiển thị tình trạng từng chỗ theo khung giờ; công cụ vẽ sơ đồ cho quản lý chi nhánh \\
\hline
2 & Dễ đặt trùng khi nhiều người đặt cùng lúc & Khóa tư vấn \mt{pg\_advisory\_xact\_lock} ở tầng dịch vụ và ràng buộc \mt{EXCLUDE} ở cơ sở dữ liệu \\
\hline
3 & Đối soát chuyển khoản thủ công & Mã VietQR có sẵn số tiền và nội dung qua PayOS; webhook xác nhận đơn tự động \\
\hline
4 & Chính sách hủy thiếu minh bạch, tính tay & Tính số tiền hoàn theo bậc thời gian; chính sách chi nhánh ưu tiên hơn chính sách chung \\
\hline
5 & Thiếu gắn kết giữa các thành viên & Hồ sơ năng lực, gợi ý đối tác, lời mời kết nối, bảng tin cộng đồng \\
\hline
6 & Tra cứu chính sách và đặt chỗ còn nhiều bước & Trợ lý ảo Gemini tra cứu dữ liệu thật và đề xuất thao tác, người dùng xác nhận trước khi thực hiện \\
\hline
\end{tabularx}
\end{table}

\section{Đề xuất hệ thống CoSpace}
\label{sec:dexuat}

\subsection{Định hướng}
CoSpace là ứng dụng web quản lý chuỗi văn phòng chia sẻ, hướng đến các chuỗi quy mô vừa và nhỏ tại Việt Nam. Hệ thống hỗ trợ cả kênh trực tuyến (đặt chỗ, thanh toán VietQR, ví MoMo) lẫn kênh tại quầy (khách vãng lai, tiền mặt), cho phép mỗi chi nhánh có bảng giá và chính sách riêng, và bổ sung các chức năng cộng đồng.

\subsection{Nhóm người dùng}
\begin{enumerate}
    \item \textbf{Khách hàng:} tìm và chọn chỗ trên sơ đồ, đặt chỗ, thanh toán, nhận mã check-in, kết nối với thành viên khác.
    \item \textbf{Nhân viên quầy:} check-in, check-out, tạo đơn cho khách vãng lai, thu tiền mặt, ghi nhận dịch vụ phát sinh.
    \item \textbf{Quản lý chi nhánh:} cấu hình mặt bằng, giá, chính sách hủy, xử lý hoàn tiền và theo dõi báo cáo của chi nhánh.
    \item \textbf{Quản trị viên hệ thống:} quản lý mạng lưới chi nhánh, chính sách chung, phân quyền và nhật ký kiểm toán.
\end{enumerate}

\subsection{Các phân hệ chức năng}
Hệ thống gồm tám phân hệ:
\begin{itemize}
    \item \textbf{Không gian và sơ đồ mặt bằng:} cấu trúc chi nhánh -- tầng -- không gian; công cụ vẽ sơ đồ và gắn phần tử trên sơ đồ với không gian.
    \item \textbf{Đặt chỗ và vòng đời đơn:} đặt trực tuyến hoặc tại quầy, giữ chỗ 15 phút, tính số đơn vị thời gian ở máy chủ; mọi chuyển trạng thái đi qua lớp \mt{BookingStateMachine}.
    \item \textbf{Thanh toán:} PayOS (VietQR), MoMo (sandbox), tiền mặt tại quầy; xử lý webhook lũy đẳng.
    \item \textbf{Hủy đơn và hoàn tiền:} chính sách nhiều bậc tính theo phút, ưu tiên chính sách chi nhánh; hàng chờ hoàn tiền do quản lý xử lý.
    \item \textbf{Vận hành quầy:} check-in bằng mã đặt chỗ hoặc QR, ghi nhận dịch vụ phát sinh, gia hạn, phí trả chỗ muộn, lịch bảo trì.
    \item \textbf{Thành viên và khuyến mãi:} mã giảm giá, voucher cá nhân, bốn hạng thành viên (Bronze, Silver, Gold, Platinum) xét theo chi tiêu thực tế.
    \item \textbf{Cộng đồng và gợi ý đối tác:} hồ sơ năng lực, bảng tin, gợi ý đối tác dựa trên hệ số Jaccard, lời mời kết nối.
    \item \textbf{Trợ lý ảo:} Google Gemini với cơ chế gọi hàm để tra cứu dữ liệu và đề xuất đặt chỗ, hủy chỗ, gọi thêm dịch vụ.
\end{itemize}

\section{Các bài toán kỹ thuật chính}
\label{sec:baitoan_kythuat}

\begin{enumerate}
    \item \textbf{Hiển thị và chỉnh sửa sơ đồ mặt bằng:} lưu bố cục tầng, vẽ lại trên trình duyệt ở mọi kích thước màn hình và gắn từng phần tử với một không gian trong cơ sở dữ liệu để tô màu theo tình trạng.
    \item \textbf{Chống đặt trùng khi có yêu cầu đồng thời:} bảo đảm hai đơn đang hiệu lực không bao giờ giao nhau trên cùng một không gian.
    \item \textbf{Xử lý thanh toán tự động:} kiểm tra chữ ký webhook, không ghi nhận lặp, xử lý tiền về sau khi đơn đã hết hạn giữ chỗ và trường hợp khách trả hai lần.
    \item \textbf{Chính sách hủy theo chi nhánh:} chọn đúng bậc chính sách theo phút, ưu tiên chính sách của chi nhánh và bảo đảm tiền hoàn cộng phí phạt bằng tổng tiền đơn.
    \item \textbf{Gợi ý đối tác:} tính điểm tương đồng từ kỹ năng, sở thích và chủ đề bài viết, cộng điểm khi cùng chi nhánh.
    \item \textbf{Trợ lý ảo có kiểm soát:} cho mô hình ngôn ngữ gọi các hàm tra cứu, nhưng mọi thao tác ghi dữ liệu phải qua bước xác nhận của người dùng.
\end{enumerate}
